\documentclass{bmstu}

\usepackage{amssymb}
\usepackage{longtable}
\usepackage[titles]{tocloft}
\usepackage{tabularx}
\usepackage{ragged2e} 

\newcolumntype{C}{>{\Centering\arraybackslash}X}

\renewcommand{\cftchapleader}{\cftdotfill{\cftdotsep}}
\renewcommand{\cftchapfont}{\normalfont}
\renewcommand{\cftchappagefont}{\normalfont}
\setlength{\cftbeforechapskip}{0pt}
\addbibresource{bibliography.bib}

\begin{document}

\maketechnicalreport
    {Информатика и системы управления}                          % #1 факультет
    {Программное обеспечение ЭВМ и информационные технологии}   % #2 кафедра
    {\textbf{Портал бронирования транспортных услуг через партнёрские сервисы}}                % #3 название
    {}                                            % #4 шифр (или {} если не нужен)
    {Е.~К.~Критинина}                                 % #5 студент
    {ИУ7-22М}                                     % #6 группа
    {Т.~Н.~Романова}                              % #7 руководитель
    {}                                            % #8 нормоконтролер (или {})


\newpage
\chapter*{ВВЕДЕНИЕ}
\addcontentsline{toc}{chapter}{ВВЕДЕНИЕ}

По данным аналитического агентства «АВТОСТАТ», на начало 2024 года рынок краткосрочной аренды и каршеринга в России насчитывает более 2,5~млн транспортных средств, распределённых между несколькими крупными операторами (Яндекс.Прокат, Делимобиль, BelkaCar и другими)~\cite{autostat2024}.

Корпоративные клиенты (компании, нуждающиеся в организации служебного транспорта для сотрудников) сталкиваются со значительными сложностями при управлении расходами на бронирование. Сотрудник, отправляющийся в командировку, вынужден самостоятельно искать доступные автомобили у различных операторов через их отдельные сайты и приложения, сравнивать цены, согласовывать поездку с руководителем по электронной почте. Отсутствует единое место для контроля расходов, управления квотами и аналитики использования транспорта.

По оценкам экспертов журнала «Автопарк», до~40\% корпоративных расходов на транспорт связаны с неэффективной организацией процесса бронирования: дублирующиеся поездки, превышение установленных бюджетов, отсутствие централизованного контроля и аудита~\cite{autopark2023}. При этом внедрение специализированной платформы управления позволяет сократить эти издержки на~15--25\% и повысить уровень контроля над командировочными расходами.

Существующие решения на рынке либо представляют собой интерфейсы отдельных операторов каршеринга (фокус на их продукты), либо являются универсальными системами управления расходами без специализированного функционала для транспорта. Отсутствуют платформы, интегрирующие API нескольких независимых операторов и предоставляющие корпоративному клиенту единый портал с поддержкой согласования заявок, управления квотами и аналитики.

Учитывая вышеперечисленные проблемы, актуальной становится разработка распределённой платформы, которая объединяет предложения от партнёрских сервисов каршеринга и проката, предоставляет единый интерфейс для бронирования, реализует бизнес-процессы согласования и контроля расходов, обеспечивает аналитику для оптимизации корпоративных транспортных расходов.

Данное техническое задание составлено для разработки распределённой системы бронирования транспортных услуг через партнёрские сервисы. Техническое задание выполнено на основе ГОСТ~34.602--89 «Информационная технология. Комплекс стандартов на автоматизированные системы. Техническое задание на создание автоматизированной системы» и ГОСТ~Р~56939--2016 «Защита информации. Разработка безопасного программного обеспечения. Общие требования».


\newpage
\chapter*{ГЛОССАРИЙ}
\addcontentsline{toc}{chapter}{ГЛОССАРИЙ}


\begin{description}

\item Партнёр -- внешний сервис (компания) предоставляющий услуги по аренде автомобилей или каршеринга, интегрированный с платформой через API
\item Бронирование (заказ) -- заявка пользователя на использование автомобиля в определённый период времени, создаваемая в системе платформы и передаваемая партнёру для исполнения
\item Внешнее бронирование -- бронирование, созданное и хранящееся в системе партнёра, идентифицируемое внешним ID
\item Заявка -- электронный документ, содержащий информацию о запрашиваемом автомобиле, сроках использования и цели поездки, требующий согласования с руководителем
\item Квота -- лимит на количество одновременных бронирований или сумму расходов, установленный для отдела или категории сотрудников
\item Микросервис -- автономный компонент системы, отвечающий за ограниченный набор функций и взаимодействующий с другими компонентами через API
\item API Gateway -- единая точка входа в систему, обеспечивающая маршрутизацию запросов, аутентификацию и авторизацию
\item Распределённая система -- система, компоненты которой размещены на различных узлах сети и взаимодействуют посредством обмена сообщениями
\item SSO (Single Sign-On) -- технология единого входа, позволяющая пользователю получить доступ к нескольким системам после однократной аутентификации
\item OAuth2 -- открытый протокол авторизации, позволяющий предоставить доступ к ресурсам без передачи логина и пароля
\item JWT (JSON Web Token) -- стандарт токена доступа, используемый для передачи данных аутентификации между сервисами
\item Circuit Breaker -- паттерн отказоустойчивости, предотвращающий каскадные сбои при недоступности одного из сервисов
\item Saga Pattern -- паттерн управления распределёнными транзакциями и компенсирующими операциями при сбое

\end{description}

\newpage
\begin{center}
    \textbf{Основания для разработки}
\end{center}

Разработка ведётся в рамках выполнения лабораторных работ по курсу
«Методология программной инженерии» на кафедре «Программное обеспечение ЭВМ и информационные технологии» факультета «Информатика и
системы управления» МГТУ им. Н.Э. Баумана.

\begin{center} 
    \textbf{Назначение разработки}
\end{center}

Разрабатываемая система должна предоставлять корпоративным клиентам и их сотрудникам единую платформу для бронирования транспортных услуг у различных партнёрских операторов каршеринга и проката. Система должна обеспечивать поиск и сравнение предложений от партнёров, автоматизацию процесса согласования поездок, управление корпоративными квотами и расходами, а также аналитику для оптимизации использования транспорта.

\begin{center} 
    \textbf{Существующие аналоги}
\end{center}

На сегодняшний день решения для управления корпоративными расходами на транспорт представлены следующими основными системами:

\begin{itemize}
    \item 1С: Управление автотранспортом~--- \url{https://solutions.1c.ru/catalog/transport};
    \item Triplist~--- \url{https://triplist.ru/};
    \item Yandex.Drive for Business~--- \url{https://business.yandex.ru/products/drive/};
    \item Delimobil Business~--- \url{https://business.delimobil.ru/};
    \item BelkaCar for Business~--- \url{https://belkacar.ru/business}.
\end{itemize}

У вышеперечисленных систем есть определённые недостатки:

\begin{itemize}
    \item закрытость на одного оператора каршеринга (Yandex.Drive, Delimobil, BelkaCar) — невозможна интеграция с альтернативными сервисами;
    \item отсутствие единого интерфейса для поиска и сравнения предложений от нескольких партнёров одновременно;
    \item слабая поддержка согласования заявок с руководителем через систему — процесс осуществляется вне платформы;
    \item ограниченные возможности по управлению корпоративными квотами и лимитами расходов;
    \item отсутствие аналитики и рекомендаций по оптимизации расходов и использования транспорта;
    \item сложность интеграции и кастомизации под специфические бизнес-процессы организации;
    \item отсутствие механизма для работы при недоступности одного из партнёров.
\end{itemize}

По сравнению с существующими аналогами разрабатываемая система должна иметь следующие преимущества:

\begin{itemize}
    \item Единый портал интеграции -- сотрудники работают с одной системой, интегрирующей несколько партнёрских сервисов;
    \item Поиск и сравнение -- система параллельно запрашивает доступность и цены у всех партнёров;
    \item Автоматизированный процесс согласования -- заявки направляются руководителю, статусы отслеживаются в системе;
    \item Управление квотами и бюджетами -- гибкая настройка лимитов по отделам и пользователям;
    \item Аналитика расходов -- консолидированные отчёты, рекомендации;
    \item Отказоустойчивость -- при недоступности одного партнёра система продолжает работать с остальными;
    \item Единый биллинг -- корпоративный клиент видит консолидированный счёт по всем партнёрам.
\end{itemize}

\begin{center}
    \textbf{Описание системы}
\end{center}

Разрабатываемый сервис должен представлять собой распределённую платформу для агрегации и оркестрации бронирований транспортных услуг от партнёрских операторов.

\textbf{Основной процесс:}

Сотрудник авторизуется в системе, используя корпоративные учётные данные (SSO). Для подачи заявки сотрудник указывает период использования, желаемый класс автомобиля и цель поездки. Заявка направляется на согласование руководителю отдела.

После одобрения заявки система запускает процесс оркестрации бронирования: параллельно отправляются запросы к API всех доступных партнёров. Партнёры возвращают информацию о доступных автомобилях, ценах и условиях. Система выбирает оптимального партнёра (по цене, рейтингу, времени доступности) и создаёт бронирование в его системе через API. Бронирование подтверждается, сотруднику отправляется уведомление с деталями автомобиля и пункта выдачи.

Система логирует все события, собирает данные о расходах и использовании для последующей аналитики. Администратор имеет доступ к консолидированным отчётам и рекомендациям по оптимизации расходов.

\textbf{Ключевые особенности:}

Система не хранит сами бронирования (они находятся в системах партнёров), а хранит заказы — внутренние объекты, содержащие ссылку на бронирование у партнёра, статус оплаты, информацию о сотруднике и параметры поездки. При недоступности партнёра система может переключиться на альтернативного поставщика или отложить бронирование в очередь.

Для неавторизованных пользователей доступна только общая информация о системе и функция входа.


\newpage

\begin{center}
    \textbf{Общие требования к системе}
\end{center}

Разрабатываемая система должна отвечать следующим требованиям.

\begin{enumerate}
    \item Разрабатываемое ПО должно поддерживать функционирование системы в режиме 24~часов, 7~дней в неделю, 365~дней в году (24/7/365) со среднегодовым временем доступности не менее 99,9\%. Допустимое время, в течение которого система недоступна, за год должно составлять $24 \cdot 365 \cdot 0{,}001 = 8{,}76$~ч.
    
    \item Время восстановления системы после сбоя не должно превышать 15~минут.
    
    \item Система должна автоматически восстанавливаться после сбоя и продолжать работу при частичной недоступности партнёров.
    
    \item Система должна поддерживать возможность «горячего» подключения новых партнёров во время работы системы без рестарта.
    
    \item Обеспечить безопасность и работоспособность за счёт отказоустойчивости и механизма Circuit Breaker при взаимодействии с партнёрскими API.
\end{enumerate}

\begin{center}
    \textbf{Требования к функциональным характеристикам}
\end{center}

\begin{enumerate}
    \item По результатам работы модуля сбора статистики медиана времени отклика системы на запросы пользователя на получение информации (просмотр доступных автомобилей, история бронирований) не должна превышать 3~секунд.
    
    \item По результатам работы модуля сбора статистики медиана времени отклика системы на запросы, добавляющие или изменяющие информацию (создание заявки, одобрение заявки, подтверждение бронирования) не должна превышать 7~секунд.
    
    \item Медиана времени отклика системы на действия пользователя должна быть менее 0,8~секунды при условии работы на рекомендованной аппаратной конфигурации, задержках между взаимодействующими сервисами менее 0,2~секунды и одновременном числе работающих пользователей менее 100 на каждый сервер.
    
    \item Система должна обеспечивать возможность запуска в современных браузерах: не менее 85\% пользователей Интернета должны пользоваться ей без какой-либо деградации функционала.
\end{enumerate}

\begin{center}
    \textbf{Функциональные требования к системе с точки зрения пользователя}
\end{center}

Система должна обеспечивать реализацию нижеуказанных функций.

\begin{enumerate}
    \item Аутентификация пользователей через корпоративный SSO.
    
    \item Авторизация пользователей с проверкой прав доступа на основе ролевой модели.
    
    \item Разделение всех пользователей на следующие роли:
    \begin{itemize}
        \item Сотрудник (авторизованный пользователь);
        \item Руководитель отдела;
        \item Администратор системы.
    \end{itemize}
    
    \item Предоставление возможностей Сотруднику, Руководителю отдела, Администратору, представленных в таблице~\ref{tab:user_roles}.
\end{enumerate}

\begin{longtable}{|p{0.25\textwidth}|p{0.7\textwidth}|}
\caption{Функции пользователей} \label{tab:user_roles} \\
\hline
\textbf{Роль} & \textbf{Функции} \\
\hline
\endfirsthead

\multicolumn{2}{c}{\tablename\ \thetable\ -- Продолжение} \\
\hline
\textbf{Роль} & \textbf{Функции} \\
\hline
\endhead

\hline
\multicolumn{2}{r}{Продолжение на следующей странице} \\
\endfoot

\hline
\endlastfoot

Сотрудник & 
\begin{enumerate}
    \item Просмотр списка доступных автомобилей от всех партнёров (поиск и фильтрация по классу, периоду, цене);
    \item Подача заявки на бронирование;
    \item Просмотр всех своих заявок и заказов с указанием статуса;
    \item Получение детальной информации по конкретной заявке/заказу;
    \item Отмена собственной заявки (до одобрения);
    \item Получение и изменение информации своего профиля.
\end{enumerate} \\
\hline

Руководитель отдела & 
\begin{enumerate}
    \item Все функции Сотрудника;
    \item Просмотр заявок сотрудников своего отдела;
    \item Одобрение или отклонение заявок сотрудников своего отдела;
    \item Просмотр статистики использования услуг бронирования сотрудниками отдела;
    \item Просмотр информации о квотах отдела.
\end{enumerate} \\
\hline

Администратор & 
\begin{enumerate}
    \item Управление партнёрами (добавление, редактирование, отключение);
    \item Управление пользователями системы (добавление, редактирование);
    \item Назначение ролей пользователям;
    \item Просмотр всех заявок и заказов;
    \item Отмена заказа (при чрезвычайных ситуациях);
    \item Управление квотами по отделам;
    \item Просмотр общей аналитики по расходам и использованию;
    \item Просмотр логов и статистики работы системы;
    \item Просмотр статуса и здоровья интеграции с партнёрами.
\end{enumerate} \\
\hline

\end{longtable}

\begin{center}
    \textbf{Входные данные}
\end{center}

Входные параметры системы представлены в таблице~\ref{tab:input_data}.

\begin{longtable}{|p{0.25\textwidth}|p{0.7\textwidth}|}
\caption{Входные данные} \label{tab:input_data} \\
\hline
\textbf{Сущность} & \textbf{Входные данные} \\
\hline
\endfirsthead

\multicolumn{2}{c}{\tablename\ \thetable\ -- Продолжение} \\
\hline
\textbf{Сущность} & \textbf{Входные данные} \\
\hline
\endhead

\hline
\multicolumn{2}{r}{Продолжение на следующей странице} \\
\endfoot

\hline
\endlastfoot

Пользователь (все роли) & 
\begin{enumerate}
    \item Фамилия, имя, отчество~--- не более 256 символов каждое поле;
    \item Логин~--- не менее 6 символов и не более 128;
    \item Пароль~--- не менее 8 символов и не более 128, минимум одна заглавная и одна строчная буква, одна цифра;
    \item Номер телефона;
    \item Электронная почта;
    \item Отдел (для сотрудников).
\end{enumerate} \\
\hline

Партнёр & 
\begin{enumerate}
    \item Идентификатор;
    \item Название компании~--- не более 256 символов;
    \item Города присутствия;
    \item Тип интеграции (REST API, другое);
    \item URL API endpoint~--- не более 512 символов;
    \item Ключи аутентификации (API key, OAuth токен);
    \item Комиссия платформы (процент);
    \item Приоритет (для выбора партнёра);
    \item Статус (активен, на тестировании, отключён);
    \item Описание~--- не более 2048 символов.
\end{enumerate} \\
\hline

Заявка на бронирование & 
\begin{enumerate}
    \item Идентификатор;
    \item Идентификатор сотрудника (инициатор заявки);
    \item Желаемый класс автомобиля;
    \item Дата и время начала использования;
    \item Дата и время окончания использования;
    \item Цель поездки~--- не более 512 символов;
    \item Статус заявки (ожидает одобрения, одобрена, отклонена, отменена);
    \item Комментарий руководителя~--- не более 512 символов;
    \item Идентификатор заказа (ссылка на созданный заказ).
\end{enumerate} \\
\hline

Заказ на бронирование & 
\begin{enumerate}
    \item Идентификатор в системе платформы (order\_id);
    \item Идентификатор заявки;
    \item Идентификатор сотрудника;
    \item Идентификатор выбранного партнёра;
    \item Внешний ID бронирования у партнёра (external\_booking\_id);
    \item Класс, период использования;
    \item Предварительная стоимость;
    \item Финальная стоимость (если известна);
    \item Статус заказа (created, partner\_selected, payment\_pending, sent\_to\_partner, confirmed, failed, cancelled, completed);
    \item Время создания и последнего обновления;
    \item Параметры платежа и информация об оплате.
\end{enumerate} \\
\hline

Квота отдела & 
\begin{enumerate}
    \item Идентификатор отдела;
    \item Максимальное количество одновременных бронирований;
    \item Максимальная сумма расходов за период;
    \item Период действия квоты (месяц/квартал/год);
    \item Текущее использование.
\end{enumerate} \\
\hline

\end{longtable}

\begin{center}
\textbf{Выходные параметры}
\end{center}

Выходными параметрами системы являются web-страницы. В зависимости от запроса и текущей роли пользователя они содержат следующую информацию (таблица~\ref{tab:output_data}).

\begin{longtable}{|p{0.25\textwidth}|p{0.7\textwidth}|}
\caption{Выходные параметры} \label{tab:output_data} \\
\hline
\textbf{Роль} & \textbf{Выходные данные} \\
\hline
\endfirsthead

\multicolumn{2}{c}{\tablename\ \thetable\ -- Продолжение} \\
\hline
\textbf{Роль} & \textbf{Выходные данные} \\
\hline
\endhead

\hline
\multicolumn{2}{r}{Продолжение на следующей странице} \\
\endfoot

\hline
\endlastfoot

Сотрудник & 
\begin{enumerate}
    \item Агрегированный список доступных автомобилей от всех партнёров:
    \begin{itemize}
        \item партнёр (название);
        \item марка и модель;
        \item класс автомобиля;
        \item технические характеристики;
        \item цена;
        \item рейтинг партнёра;
        \item время доставки.
    \end{itemize}
    \item Детальная информация о пользователе, вошедшем в систему;
    \item Список заявок пользователя с указанием статуса;
    \item Детальная информация о заявке:
    \begin{itemize}
        \item параметры бронирования;
        \item статус одобрения;
        \item информация о заказе (если создан);
        \item история изменений.
    \end{itemize}
    \item История заказов с информацией о затратах.
\end{enumerate} \\
\hline

Руководитель отдела & 
\begin{enumerate}
    \item Все выходные данные роли «Сотрудник»;
    \item Список заявок сотрудников своего отдела с указанием:
    \begin{itemize}
        \item ФИО сотрудника;
        \item параметры запрашиваемой услуги;
        \item статус заявки;
        \item приблизительная стоимость.
    \end{itemize}
    \item Статистика отдела:
    \begin{itemize}
        \item количество одобренных бронирований;
        \item использованная квота;
        \item общие расходы за период;
        \item тренды использования.
    \end{itemize}
\end{enumerate} \\
\hline

Администратор & 
\begin{enumerate}
    \item Список всех партнёров с информацией:
    \begin{itemize}
        \item название, города присутствия;
        \item статус интеграции и здоровье API;
        \item последнее время обновления;
        \item количество успешных/неудачных запросов.
    \end{itemize}
    \item Список всех пользователей системы;
    \item Полный реестр заявок и заказов;
    \item Консолидированная аналитика:
    \begin{itemize}
        \item общие расходы по отделам;
        \item распределение использования по партнёрам;
        \item популярные направления и время бронирования;
        \item анализ стоимости (средняя, мин, макс);
        \item статистика одобрений/отклонений заявок.
    \end{itemize}
    \item Рекомендации:
    \begin{itemize}
        \item выявление избыточных расходов;
        \item анализ эффективности партнёров;
        \item предложения по оптимизации.
    \end{itemize}
    \item Логи работы системы и метрики производительности.
\end{enumerate} \\
\hline

\end{longtable}

\begin{center}
    \textbf{Топология системы}
\end{center}

На рисунке~\ref{fig:topology} изображена топология разрабатываемой распределённой системы.

\begin{figure}[h]
    \centering
    % \includegraphics[width=0.95\textwidth]{topology.png}
    \caption{Топология распределённой системы бронирования транспортных услуг}
    \label{fig:topology}
\end{figure}

Система будет состоять из фронтенда и следующих микросервисов:

\begin{itemize}
    \item API Gateway;
    \item сервис аутентификации и авторизации;
    \item сервис оркестрации бронирований;
    \item сервис аналитики;
    \item сервис статистики и логирования.
\end{itemize}

\textbf{Фронтенд} принимает запросы от пользователей по протоколу HTTP и анализирует их. На основе проведённого анализа выполняет запросы к микросервисам бэкенда, агрегирует ответы и отсылает их пользователю.

\textbf{API Gateway}~--- единая точка входа в систему. Все запросы пользователей маршрутизируются через данный сервис. API Gateway отвечает за валидацию входных данных, rate limiting, и трансляцию запросов к соответствующим микросервисам.

\textbf{Сервис аутентификации и авторизации} отвечает за:
\begin{itemize}
    \item аутентификацию пользователя через корпоративный SSO или OAuth2;
    \item авторизацию пользователя на основе ролевой модели;
    \item выдачу и управление JWT-токенами;
    \item управление сессиями пользователей;
    \item проверку прав доступа к API endpoints.
\end{itemize}

\textbf{Сервис оркестрации бронирований (Booking Orchestrator Service)} реализует основную бизнес-логику системы:
\begin{itemize}
    \item приём и обработку заявок на бронирование от пользователей;
    \item проверку квот и лимитов расходов;
    \item управление процессом согласования заявок с руководителями;
    \item выбор оптимального партнёра на основе критериев (цена, доступность, рейтинг, SLA);
    \item взаимодействие с API партнёрских сервисов для создания бронирований;
    \item обработку платежей (авторизация, захват, возврат средств);
    \item управление жизненным циклом заказа (создание, подтверждение, отмена, завершение);
    \item обработку ошибок и реализацию компенсирующих транзакций (Saga pattern);
    \item синхронизацию статусов с партнёрами;
    \item логирование всех операций для аудита.
\end{itemize}

\textbf{Сервис аналитики} реализует следующие функции:
\begin{itemize}
    \item агрегацию данных о расходах по отделам и пользователям;
    \item анализ использования услуг и популярных направлений;
    \item мониторинг эффективности партнёров;
    \item расчёт метрик производительности;
    \item формирование отчётов и рекомендаций;
    \item выявление аномалий и неэффективного использования;
    \item прогнозирование спроса.
\end{itemize}

\textbf{Сервис статистики и логирования} отвечает за:
\begin{itemize}
    \item сбор событий от всех сервисов системы;
    \item сохранение логов в централизованное хранилище;
    \item формирование метрик производительности (latency, throughput, error rate);
    \item мониторинг здоровья интеграции с партнёрами;
    \item предоставление данных для дашбордов мониторинга;
    \item архивирование логов.
\end{itemize}

\textbf{Партнёрские адаптеры} (Partner Adapters) — набор модулей для интеграции с различными партнёрскими сервисами:
\begin{itemize}
    \item реализация API-интеграции с каждым партнёром;
    \item нормализация данных партнёра в единый формат платформы;
    \item обработка ошибок и таймаутов;
    \item реализация паттерна Circuit Breaker;
    \item кеширование информации о доступных автомобилях и тарифах;
    \item получение вебхуков от партнёров об изменении статусов.
\end{itemize}

Каждый партнёр может быть подключен один или несколько раз (для разных API версий или конфигураций).

\begin{center}
    \textbf{Процесс бронирования}
\end{center}

\textbf{Этап 1: Подача заявки}

Сотрудник заполняет форму с параметрами требуемой услуги (период, класс автомобиля, цель поездки) и отправляет заявку. Сервис оркестрации создаёт объект ApplicationRequest и переводит его в статус «pending\_approval».

\textbf{Этап 2: Согласование}

Заявка направляется руководителю отдела на согласование. Руководитель проверяет, укладывается ли запрос в квоту отдела и соответствует ли он политикам компании. При одобрении заявка переходит в статус «approved». При отклонении — в статус «rejected».

\textbf{Этап 3: Поиск и выбор партнёра}

При одобрении заявки сервис оркестрации:
\begin{enumerate}
    \item определяет доступных партнёров;
    \item отправляет параллельные запросы ко всем партнёрам для получения информации о доступных автомобилях и ценах;
    \item агрегирует результаты;
    \item применяет алгоритм выбора партнёра (оптимизация по цене, рейтингу, времени доступности);
    \item выбирает лучшего партнёра.
\end{enumerate}

\textbf{Этап 4: Обработка платежа}

Если платежи проходят через платформу:
\begin{enumerate}
    \item система инициирует авторизацию суммы на корпоративной карте/счёте;
    \item если авторизация прошла — переходит к следующему этапу;
    \item если авторизация отклонена — заказ переводится в статус «payment\_failed» и уведомляет пользователя.
\end{enumerate}

\textbf{Этап 5: Создание бронирования у партнёра}

Система отправляет запрос на создание бронирования в API выбранного партнёра с параметрами заказа. Партнёр возвращает ID бронирования (external\_booking\_id) и деталей (марка автомобиля, место выдачи, условия и т.д.). Система сохраняет внешний ID и переводит заказ в статус «confirmed».

\textbf{Этап 6: Захват платежа и уведомление}

\begin{enumerate}
    \item если платёж был авторизован — система захватывает средства (capture);
    \item система отправляет уведомление сотруднику с деталями бронирования и инструкциями;
    \item информация логируется для аудита и аналитики.
\end{enumerate}

\textbf{Обработка ошибок:}

\begin{itemize}
    \item Если партнёр вернул ошибку при создании бронирования — система пытается перебронировать у альтернативного партнёра (если доступен).
    \item Если все партнёры недоступны — система помещает заказ в очередь Kafka для повторной попытки.
    \item Если платёж был захвачен, но бронирование не удалось — система выполняет возврат средств (refund) и логирует инцидент.
    \item В случае критической ошибки — система уведомляет администратора и пользователя.
\end{itemize}

\begin{center}
    \textbf{Требования к надёжности и отказоустойчивости}
\end{center}

Разрабатываемая система должна обеспечивать непрерывную работу даже при выходе из строя отдельных компонентов инфраструктуры. Для достижения требуемого уровня надёжности необходимо выполнение следующих условий:

\begin{enumerate}
    \item Архитектура системы должна базироваться на принципах сервис-ориентированного подхода с выделением независимых микросервисов, каждый из которых отвечает за отдельную область бизнес-логики.
    
    \item Инфраструктура должна предусматривать распределённое развёртывание сервисов на нескольких физических или виртуальных серверах с поддержкой горизонтального масштабирования при росте нагрузки.
    
    \item API Gateway выполняет функции централизованной маршрутизации запросов и распределения нагрузки между экземплярами сервисов.
    
    \item Система должна поддерживать динамическое добавление новых партнёров и масштабирование существующих сервисов без остановки работы всей системы.
    
    \item Каждый микросервис должен автоматически восстанавливаться после сбоя. Максимальное время восстановления работоспособности всей системы после отказа критичного компонента не должно превышать 15~минут.
    
    \item При недоступности партнёрского API система должна продолжать функционирование в режиме деградации: запросы помещаются в очередь для повторной попытки, система переключается на альтернативных партнёров (если доступны). Реализация паттерна Circuit Breaker обязательна для предотвращения каскадных отказов.
    
    \item Для управления распределёнными транзакциями при взаимодействии с партнёрами необходимо использовать паттерн Saga. При сбое на любом этапе процесса (например, платёж прошёл, но бронирование у партнёра неудачно) система должна выполнять компенсирующие транзакции (откат платежа, отмена частичных бронирований и т.д.).
    
    \item Для обеспечения асинхронной обработки и eventual consistency необходимо использовать брокер сообщений Apache Kafka. При сбое обработки события оно должно сохраняться в очереди для повторной попытки с экспоненциальной задержкой.
    
    \item Резервное копирование баз данных должно выполняться автоматически не реже одного раза в сутки с сохранением копий на отдельном хранилище. Время восстановления данных из резервной копии не должно превышать 2~часов.
    
    \item Система должна поддерживать идемпотентность операций для безопасного повтора запросов при сбоях без создания дубликатов.
\end{enumerate}

\begin{center}
    \textbf{Требования к производительности и временным характеристикам}
\end{center}

Система должна обеспечивать высокую скорость обработки запросов пользователей даже при пиковых нагрузках. Для контроля показателей производительности используется сервис статистики, который собирает метрики в режиме реального времени. Требуемые характеристики представлены в таблице~\ref{tab:performance}.

\begin{table}[h]
\centering
\caption{Требования к производительности системы}
\label{tab:performance}
\begin{tabularx}{\textwidth}{|X|c|}
\hline
\textbf{Показатель производительности} & \textbf{Требуемое значение} \\
\hline
Медиана времени отклика на запросы получения информации (поиск автомобилей, история заказов) & $\leq 3$~с \\
\hline
Медиана времени отклика на запросы изменения данных (создание заявки, одобрение, запуск бронирования) & $\leq 7$~с \\
\hline
Медиана времени отклика на пользовательские действия при штатной нагрузке & $\leq 0{,}8$~с \\
\hline
Задержка при межсервисном взаимодействии (latency между микросервисами) & $\leq 0{,}2$~с \\
\hline
Медиана времени отклика при запросах к API партнёров & $\leq 2$~с \\
\hline
Годовая доступность системы (uptime) & $\geq 99{,}9$~\% \\
\hline
Допустимое суммарное время недоступности системы в течение года & $\leq 8{,}76$~ч \\
\hline
\end{tabularx}
\end{table}

Дополнительные требования к производительности:

\begin{enumerate}
    \item При временной недоступности партнёра запросы на бронирование должны быть помещены в очередь Kafka для повторной попытки после восстановления доступности.
    
    \item Система должна поддерживать автоматическое горизонтальное масштабирование (Horizontal Pod Autoscaler в Kubernetes) при превышении пороговых значений нагрузки: CPU $>$ 70\%, Memory $>$ 80\%.
    
    \item Пиковая нагрузка не должна приводить к деградации производительности более чем на 20\% от указанных значений медианного времени отклика.
    
    \item Кеширование данных о доступных автомобилях и тарифах от партнёров должно сократить время поиска до менее 0,5~секунды при условии попадания в кеш.
\end{enumerate}

\begin{center}
    \textbf{Требования к эксплуатации и техническому обеспечению}
\end{center}

\begin{center}
    \textbf{Условия эксплуатации серверного оборудования}
\end{center}

Серверы, на которых развёртывается система, должны эксплуатироваться в условиях, соответствующих следующим требованиям:

\begin{enumerate}
    \item Температурный режим в серверном помещении должен поддерживаться в диапазоне от $+1$\textdegree C до $+40$\textdegree C.
    
    \item Относительная влажность воздуха не должна превышать 80\% при температуре $+25$\textdegree C.
    
    \item Все серверы должны быть подключены к источникам бесперебойного питания (ИБП), обеспечивающим автономную работу оборудования не менее 30~минут при полной нагрузке.
    
    \item Плановое техническое обслуживание должно проводиться не чаще одного раза в месяц в нерабочее время.
    
    \item Обновление программного обеспечения системы должно выполняться без остановки её работы. Для этого используется стратегия rolling update в Kubernetes.
\end{enumerate}

\begin{center}
    \textbf{Минимальные требования к серверному оборудованию}
\end{center}

\begin{table}[h]
\centering
\caption{Требования к серверному оборудованию}
\label{tab:server_requirements}
\begin{tabularx}{\textwidth}{|l|X|c|c|}
\hline
\textbf{Тип сервера} & \textbf{Роль в системе} & \textbf{CPU} & \textbf{RAM} \\
\hline
Вычислительный узел & API Gateway, Auth, Analytics & $\geq 2$~ГГц, 2~ядра & $\geq 4$~ГБ \\
\hline
Узел с хранилищем & Booking Orchestrator, Statistics & $\geq 2{,}5$~ГГц, 4~ядра & $\geq 8$~ГБ \\
\hline
Брокер сообщений & Apache Kafka & $\geq 2{,}5$~ГГц, 4~ядра & $\geq 16$~ГБ \\
\hline
\end{tabularx}
\end{table}

\begin{table}[h]
\centering
\caption{Требования к дисковому пространству и СУБД}
\label{tab:storage_requirements}
\begin{tabularx}{\textwidth}{|l|X|X|}
\hline
\textbf{Тип сервера} & \textbf{Дисковое пространство} & \textbf{СУБД} \\
\hline
Вычислительный узел & ОС $\geq 20$~ГБ, приложение $\geq 500$~МБ & не требуется \\
\hline
Узел с хранилищем & ОС $\geq 20$~ГБ, приложение $\geq 500$~МБ, БД $\geq 500$~ГБ & PostgreSQL 14+ \\
\hline
Брокер сообщений & ОС $\geq 20$~ГБ, Kafka $\geq 1$~ТБ (логи событий) & не требуется \\
\hline
\end{tabularx}
\end{table}

Дополнительные требования:

\begin{enumerate}
    \item Операционная система сервера: Linux Ubuntu 22.04 LTS или более поздняя стабильная версия.
    
    \item Для продуктивной среды рекомендуется использовать SSD-диски для баз данных и хранилища Kafka.
    
    \item Серверы должны быть объединены в локальную сеть с пропускной способностью не менее 1~Гбит/с.
\end{enumerate}

\begin{center}
    \textbf{Требования к клиентским устройствам}
\end{center}

\begin{enumerate}
    \item Поддерживаемые браузеры: последние стабильные версии Google Chrome, Mozilla Firefox, Safari, Microsoft Edge.
    
    \item Минимальное разрешение экрана: $1280 \times 720$ (HD).
    
    \item Интерфейс должен быть адаптивным (responsive design).
    
    \item Сетевое подключение: Ethernet 1000BASE-T или Wi-Fi стандарта 802.11n и выше. Минимальная скорость: 5~Мбит/с.
    
    \item Поддержка JavaScript и cookies в браузере (обязательно).
\end{enumerate}

\begin{center}
    \textbf{Требования к информационной и программной совместимости}
\end{center}

\begin{center}
    \textbf{Протоколы взаимодействия и форматы данных}
\end{center}

\begin{enumerate}
    \item Все коммуникации между микросервисами реализуются по протоколу HTTP/HTTPS с использованием RESTful API. Формат обмена данными~--- JSON.
    
    \item Web-интерфейс формирует HTML-страницы в ответ на HTTP-запросы. Для динамического контента используется JavaScript с fetch API.
    
    \item Доступ к базам данных разрешён исключительно микросервисам-владельцам соответствующих схем данных. Прямые межбазовые запросы запрещены.
    
    \item Асинхронный обмен событиями между сервисами осуществляется через брокер сообщений Apache Kafka. Формат сообщений~--- JSON с полями: \texttt{event\_type}, \texttt{timestamp}, \texttt{source\_service}, \texttt{payload}.
    
    \item Аутентификация и авторизация реализуются с использованием протокола OAuth2 и токенов JWT.
    
    \item Интеграция с партнёрскими API выполняется через REST с использованием API ключей или OAuth2.
\end{enumerate}

\begin{center}
    \textbf{Технологический стек}
\end{center}

Для разработки системы рекомендуется использовать следующий набор технологий:

\begin{itemize}
    \item \textbf{Язык программирования бэкенда}: Java 17 LTS или выше
    \item \textbf{Фреймворк для микросервисов}: Spring Boot 3.x с Spring Cloud
    \item \textbf{API Gateway}: Spring Cloud Gateway
    \item \textbf{СУБД}: PostgreSQL 14 или выше
    \item \textbf{ORM}: Hibernate (JPA)
    \item \textbf{Брокер сообщений}: Apache Kafka 3.x
    \item \textbf{Кеширование}: Redis
    \item \textbf{Аутентификация}: Spring Security с OAuth2 и JWT
    \item \textbf{HTTP клиент}: Spring WebClient или OkHttp для взаимодействия с API партнёров
    \item \textbf{Контейнеризация}: Docker
    \item \textbf{Оркестрация контейнеров}: Kubernetes
    \item \textbf{Сборка проекта}: Maven или Gradle
    \item \textbf{CI/CD}: GitHub Actions
    \item \textbf{Мониторинг и логирование}: Prometheus + Grafana, ELK Stack
    \item \textbf{Service Discovery}: Spring Cloud Netflix Eureka или Consul
    \item \textbf{Circuit Breaker}: Resilience4j
    \item \textbf{Фронтенд}: HTML5, CSS3, JavaScript (React или Vue.js)
\end{itemize}

\begin{center}
    \textbf{Интеграция с внешними сервисами}
\end{center}

Система должна поддерживать интеграцию со следующими категориями внешних сервисов:

\begin{enumerate}
    \item \textbf{Партнёрские API каршеринга и проката}:
    \begin{itemize}
        \item Яндекс.Прокат
        \item Делимобиль
        \item BelkaCar
        \item Локальные прокаты
        \item Другие операторы каршеринга
    \end{itemize}
    Интеграция выполняется через REST API с использованием API ключей или OAuth2.
    
    \item \textbf{Провайдеры OAuth2 для аутентификации}:
    \begin{itemize}
        \item Корпоративный SSO (Keycloak, Azure AD, Okta)
        \item Социальные сети: VK, Yandex (для демонстрационных целей)
    \end{itemize}
    
    \item \textbf{Платёжные системы}:
    \begin{itemize}
        \item Yandex.Kassa
        \item Sberbank Ecommerce
        \item Другие эквайринг-провайдеры
    \end{itemize}
    Для авторизации и захвата платежей.
    
    \item \textbf{Корпоративная почта}:
    \begin{itemize}
        \item Отправка уведомлений сотрудникам о статусе заявок и заказов через SMTP
        \item JavaMail API или Spring Mail
    \end{itemize}
\end{enumerate}

Все интеграции с внешними API должны быть изолированы в отдельные адаптеры (Adapter pattern) с использованием Circuit Breaker и retry logic для обеспечения отказоустойчивости.

\begin{center}
    \textbf{Требования к информационной безопасности}
\end{center}

\begin{enumerate}
    \item \textbf{Защита учётных данных}:
    \begin{itemize}
        \item Пароли должны храниться в виде хешей BCrypt с cost factor не менее 12;
        \item Использование устаревших алгоритмов (MD5, SHA-1) запрещено.
    \end{itemize}
    
    \item \textbf{Шифрование передаваемых данных}:
    \begin{itemize}
        \item Все соединения клиент-сервер должны использовать HTTPS с TLS~1.2 или выше;
        \item Критичные данные при передаче между микросервисами должны шифроваться.
    \end{itemize}
    
    \item \textbf{Защита от веб-уязвимостей}:
    \begin{itemize}
        \item SQL-инъекции: использование параметризованных запросов;
        \item XSS: экранирование данных перед отображением;
        \item CSRF: защита через токены;
        \item Небезопасная десериализация: валидация входных данных.
    \end{itemize}
    
    \item \textbf{Контроль доступа}:
    \begin{itemize}
        \item Ролевая модель доступа (RBAC) на уровне API Gateway;
        \item Проверка прав доступа через @PreAuthorize аннотации;
        \item Логирование всех действий администратора.
    \end{itemize}
    
    \item \textbf{Управление сессиями}:
    \begin{itemize}
        \item JWT токены с коротким сроком жизни (15--30 минут для Access Token);
        \item Refresh Token до 7 дней;
        \item Токены добавляются в чёрный список (Redis) при logout.
    \end{itemize}
    
    \item \textbf{Логирование и аудит}:
    \begin{itemize}
        \item Все значимые операции логируются с указанием user\_id, типа операции, timestamp, IP-адреса, результата;
        \item Логи хранятся не менее 1~года;
        \item Доступ к логам предоставляется только администраторам.
    \end{itemize}
    
    \item \textbf{Валидация входных данных}:
    \begin{itemize}
        \item Все входные данные проходят валидацию на API Gateway;
        \item Проверка: тип данных, диапазон значений, длина строк, формат;
        \item При ошибке возвращается HTTP 400 без раскрытия внутренней структуры системы.
    \end{itemize}
    
    \item \textbf{Rate Limiting}:
    \begin{itemize}
        \item Ограничение: не более 100 запросов в минуту от одного IP;
        \item Для критичных операций: не более 10 запросов в минуту на пользователя;
        \item При превышении возвращается HTTP 429.
    \end{itemize}
    
    \item \textbf{Управление секретами}:
    \begin{itemize}
        \item Пароли БД, API ключи, JWT ключи хранятся в Kubernetes Secrets;
        \item Запрещено хранение секретов в исходном коде или открытых переменных окружения.
    \end{itemize}
    
    \item \textbf{Обновления безопасности}:
    \begin{itemize}
        \item Регулярное обновление библиотек и фреймворков;
        \item Автоматизированная проверка зависимостей (OWASP Dependency-Check, Snyk).
    \end{itemize}
\end{enumerate}

Система должна соответствовать требованиям ГОСТ~Р~56939--2016.

\begin{center}
    \textbf{Требования к хранению данных в базе данных}
\end{center}

Данные системы хранятся в реляционных СУБД. Каждый микросервис владеет собственной схемой данных. Прямые межбазовые запросы запрещены — обмен осуществляется через REST API или Kafka.

Рекомендуемая СУБД: PostgreSQL 14+.

\begin{center}
    \textbf{Общие требования к организации данных}
\end{center}

\begin{enumerate}
    \item \textbf{Обеспечение целостности}:
    \begin{itemize}
        \item Ссылочная целостность на уровне СУБД через Foreign Key;
        \item Стратегия удаления: CASCADE или SET NULL.
    \end{itemize}
    
    \item \textbf{Транзакционность}:
    \begin{itemize}
        \item Операции выполняются в транзакциях с уровнем изоляции READ COMMITTED;
        \item При ошибке выполняется ROLLBACK;
        \item Использование @Transactional в Spring.
    \end{itemize}
    
    \item \textbf{Индексация}:
    \begin{itemize}
        \item Индексы на полях фильтрации и соединения;
        \item Обязательные индексы: первичные ключи, внешние ключи, user\_id, статусы, даты.
    \end{itemize}
    
    \item \textbf{Производительность}:
    \begin{itemize}
        \item Время SELECT-запросов $\leq 200$~мс при 100 пользователей;
        \item Аналитические запросы $\leq 3$~с.
    \end{itemize}
\end{enumerate}

\begin{center}
    \textbf{Резервное копирование и восстановление}
\end{center}

\begin{enumerate}
    \item Полное резервное копирование баз данных ежедневно в нерабочее время.
    
    \item Резервные копии хранятся на отдельном физическом носителе не менее 30~дней.
    
    \item Резервные копии должны быть зашифрованы.
    
    \item Процедура восстановления тестируется не реже одного раза в квартал.
    
    \item Время восстановления не должно превышать 2~часов.
\end{enumerate}

\begin{center}
    \textbf{Безопасность данных}
\end{center}

\begin{enumerate}
    \item \textbf{Защита паролей}:
    \begin{itemize}
        \item Хеши BCrypt с cost factor $\geq 12$.
    \end{itemize}
    
    \item \textbf{Шифрование персональных данных}:
    \begin{itemize}
        \item Рекомендуется шифрование AES-256 для ФИО, телефонов, email.
    \end{itemize}
    
    \item \textbf{Разграничение доступа}:
    \begin{itemize}
        \item Каждый микросервис подключается под отдельной учётной записью;
        \item Минимально необходимые привилегии (least privilege principle).
    \end{itemize}
    
    \item \textbf{Логирование изменений}:
    \begin{itemize}
        \item Все операции изменения данных логируются;
        \item Логи хранятся не менее 1~года.
    \end{itemize}
\end{enumerate}

\begin{center}
    \textbf{Состав сущностей по сервисам}
\end{center}

\begin{enumerate}
    \item \textbf{Сервис аутентификации} (схема \texttt{auth\_service}):
    \begin{itemize}
        \item \texttt{users}~--- пользователи;
        \item \texttt{roles}~--- роли;
        \item \texttt{user\_roles}~--- связь пользователей и ролей;
        \item \texttt{sessions}~--- активные сессии;
        \item \texttt{departments}~--- отделы.
    \end{itemize}
    
    \item \textbf{Сервис оркестрации бронирований} (схема \texttt{booking\_service}):
    \begin{itemize}
        \item \texttt{application\_requests}~--- заявки на бронирование;
        \item \texttt{orders}~--- заказы (объекты координирования);
        \item \texttt{department\_quotas}~--- квоты отделов;
        \item \texttt{payments}~--- информация о платежах;
        \item \texttt{partners}~--- конфигурация партнёров.
    \end{itemize}
    
    \item \textbf{Сервис статистики} (схема \texttt{statistics\_service}):
    \begin{itemize}
        \item \texttt{events}~--- события для мониторинга;
        \item \texttt{logs}~--- логи операций;
        \item \texttt{metrics}~--- метрики производительности.
    \end{itemize}
    
    \item \textbf{Сервис аналитики} (схема \texttt{analytics\_service}):
    \begin{itemize}
        \item \texttt{expense\_aggregates}~--- агрегированные расходы;
        \item \texttt{partner\_performance}~--- метрики партнёров;
        \item \texttt{user\_statistics}~--- статистика пользователей.
    \end{itemize}
\end{enumerate}

\begin{center}
    \textbf{Дополнительные требования}
\end{center}

\begin{enumerate}
    \item Для полей с ограниченным набором значений (статусы, роли) использовать ENUM или справочные таблицы.
    
    \item Для расширяемых данных (параметры партнёров, метаинформация) допускается использование JSONB в PostgreSQL.
    
    \item Все таблицы должны содержать поля \texttt{created\_at} и \texttt{updated\_at}.
    
    \item Логи старше 1~года должны автоматически архивироваться или удаляться.
\end{enumerate}

\begin{center}
    \textbf{Использование Kubernetes для оркестрации}
\end{center}

\begin{center}
    \textbf{Назначение Kubernetes в системе}
\end{center}

Kubernetes применяется для управления контейнерами микросервисов и обеспечивает:

\begin{enumerate}
    \item \textbf{Горизонтальную масштабируемость}~--- автоматическое увеличение количества экземпляров сервисов при росте нагрузки.
    
    \item \textbf{Автоматическое восстановление}~--- перезапуск упавших контейнеров и перенос pods на работающие узлы. Время восстановления не более 15~минут.
    
    \item \textbf{Декларативное управление}~--- конфигурация хранится в YAML-манифестах в Git.
    
    \item \textbf{Балансировку нагрузки}~--- распределение запросов между репликами сервисов.
    
    \item \textbf{Обновление без простоя}~--- Rolling Update старых версий на новые.
\end{enumerate}

\begin{center}
    \textbf{Архитектура развёртывания}
\end{center}

\begin{enumerate}
    \item \textbf{Контейнеризация}:
    \begin{itemize}
        \item Каждый микросервис в Docker-образе;
        \item Образы собираются в CI/CD pipeline (GitHub Actions);
        \item Публикуются в Container Registry.
    \end{itemize}
    
    \item \textbf{Kubernetes-ресурсы}:
    \begin{itemize}
        \item \textbf{Deployment}~--- управление pods и репликами;
        \item \textbf{Service}~--- балансировка нагрузки и обнаружение сервисов;
        \item \textbf{ConfigMap}~--- конфигурационные файлы;
        \item \textbf{Secret}~--- защищённое хранение паролей и ключей;
        \item \textbf{Ingress}~--- маршрутизация HTTPS запросов;
        \item \textbf{HorizontalPodAutoscaler}~--- автоматическое масштабирование (CPU > 70\%, Memory > 80\%).
    \end{itemize}
    
    \item \textbf{Базы данных}:
    \begin{itemize}
        \item Рекомендуется размещение PostgreSQL вне Kubernetes на выделенных серверах;
        \item Для тестовых сред допускается StatefulSet с PersistentVolumes.
    \end{itemize}
    
    \item \textbf{Apache Kafka}:
    \begin{itemize}
        \item Может быть развёрнута в Kubernetes с использованием Strimzi Operator;
        \item Минимум 3 broker для репликации;
        \item Альтернатива: управляемый сервис (Confluent Cloud, AWS MSK).
    \end{itemize}
    
    \item \textbf{Стратегия обновления}:
    \begin{itemize}
        \item Rolling Update с \texttt{maxUnavailable: 1}, \texttt{maxSurge: 1};
        \item Readiness Probe для проверки готовности;
        \item Liveness Probe для проверки работоспособности.
    \end{itemize}
    
    \item \textbf{Мониторинг}:
    \begin{itemize}
        \item Prometheus для метрик;
        \item Grafana для визуализации;
        \item ELK Stack для логирования.
    \end{itemize}
\end{enumerate}

\begin{center}
    \textbf{Пример конфигурации}
\end{center}

Каждый микросервис описывается манифестами:

\begin{itemize}
    \item \texttt{deployment.yaml}~--- Deployment с репликами, образом, переменными окружения;
    \item \texttt{service.yaml}~--- Service для балансировки;
    \item \texttt{configmap.yaml}~--- конфигурационные параметры;
    \item \texttt{secret.yaml}~--- секретные данные;
    \item \texttt{hpa.yaml}~--- автомасштабирование.
\end{itemize}

Манифесты хранятся в директории \texttt{/k8s} репозитория и применяются командой \texttt{kubectl apply -f k8s/}.

\newpage
\chapter*{ЗАКЛЮЧЕНИЕ}
\addcontentsline{toc}{chapter}{ЗАКЛЮЧЕНИЕ}

Разрабатываемая система представляет собой распределённую платформу бронирования транспортных услуг, которая решает задачу комплексного управления корпоративными расходами на каршеринг и прокат автомобилей. Интеграция с несколькими партнёрскими сервисами позволяет корпоративным клиентам получить единый интерфейс для поиска, сравнения и бронирования, с сохранением полного контроля над расходами и квотами.

Архитектура системы на основе микросервисов, использование Kafka для асинхронной обработки, реализация Saga pattern для управления распределёнными транзакциями и применение Circuit Breaker для отказоустойчивости обеспечивают надёжность и масштабируемость даже при росте нагрузки и расширении количества партнёрских интеграций.

\makebibliography

\end{document}